\chapter*{符号表}
\addcontentsline{toc}{chapter}{符号表}
\markboth{符号表}{符号表}

原书页码沿用英文版符号索引，数学字体随本译本正文；译本定位链接到对应章、节的实际起页，不表示该符号首次出现的精确页码。空白说明沿用原书未加释义的条目。\par\medskip

{\footnotesize
\begin{longtable}{@{}p{.27\linewidth}p{.31\linewidth}p{.15\linewidth}p{.17\linewidth}@{}}
\toprule 符号 & 说明 & 原书页码 & 译本定位\\\midrule
\endfirsthead
\toprule 符号 & 说明 & 原书页码 & 译本定位\\\midrule
\endhead
\(\alpha\) &  & 20 & \hyperref[sec:ch01-exercises]{第\pageref*{sec:ch01-exercises}页}\\
\(\alpha_H\) & 作用形式 & 185 & \hyperref[sec:functions-forms-covers]{第\pageref*{sec:functions-forms-covers}页}\\
\(\mathscr{A}_H\) & 作用泛函 & 158 & \hyperref[sec:action-functional]{第\pageref*{sec:action-functional}页}\\
\(\beta_k\) & 贝蒂数 & 85 & \hyperref[sec:euler-characteristic]{第\pageref*{sec:euler-characteristic}页}\\
\(c_1(TW)\) & 第一陈类 & 559 & \hyperref[sec:appB-chern-classes]{第\pageref*{sec:appB-chern-classes}页}\\
\(c_k(f)\) &  & 85 & \hyperref[sec:euler-characteristic]{第\pageref*{sec:euler-characteristic}页}\\
\(\mathcal{C}^k(\cdot,\cdot)\) &  & 544 & \hyperref[sec:appA-transversality]{第\pageref*{sec:appA-transversality}页}\\
\(\mathcal{C}^{\infty}(\cdot,\cdot)\) &  & 544 & \hyperref[sec:appA-transversality]{第\pageref*{sec:appA-transversality}页}\\
\(\mathcal{C}^{\infty}_{\mathrm{loc}}\) &  & 167 & \hyperref[sec:space-of-solutions]{第\pageref*{sec:space-of-solutions}页}\\
\(\mathcal{C}^{\infty}_{\searrow}(x,y)\) &  & 226 & \hyperref[sec:banach-manifold-p1p]{第\pageref*{sec:banach-manifold-p1p}页}\\
\(\mathcal{C}^{\infty}_{\varepsilon}\) &  & 223，233 & \hyperref[sec:results]{第\pageref*{sec:results}页}\\
\(C_k\) &  & 53，305 & \hyperref[sec:definition-of-the-complex]{第\pageref*{sec:definition-of-the-complex}页}\\
\(C(u)\) &  & 247 & \hyperref[sec:transversality]{第\pageref*{sec:transversality}页}\\
\((C_\star(f),\partial_X)\) & 莫尔斯复形 & 53 & \hyperref[sec:definition-of-the-complex]{第\pageref*{sec:definition-of-the-complex}页}\\
\((C_\star(H),\partial)\) & 弗洛尔复形 & 305 & \hyperref[sec:trajectory-spaces]{第\pageref*{sec:trajectory-spaces}页}\\
\(CM_\star(H,J)\) & 莫尔斯复形 & 359 & \hyperref[sec:floer-to-morse-results]{第\pageref*{sec:floer-to-morse-results}页}\\
\(CF_\star(H,J)\) & 弗洛尔复形 & 359 & \hyperref[sec:floer-to-morse-results]{第\pageref*{sec:floer-to-morse-results}页}\\
\(\mathrm{Crit}\) &  & 29 & \hyperref[sec:gradients-pseudo-charts]{第\pageref*{sec:gradients-pseudo-charts}页}\\
\(\chi(V)\) & V 的欧拉示性数 & 85 & \hyperref[sec:euler-characteristic]{第\pageref*{sec:euler-characteristic}页}\\
\((d^2f)\) &  & 7 & \hyperref[sec:definition-of-morse-functions]{第\pageref*{sec:definition-of-morse-functions}页}\\
\(\partial_X\) &  & 53 & \hyperref[sec:definition-of-the-complex]{第\pageref*{sec:definition-of-the-complex}页}\\
\(\overline{\partial}\) &  & 237 & \hyperref[sec:linearization-floer]{第\pageref*{sec:linearization-floer}页}\\
\(\partial_-V\) &  & 75 & \hyperref[sec:cobordisms]{第\pageref*{sec:cobordisms}页}\\
\(\partial_+V\) &  & 75 & \hyperref[sec:cobordisms]{第\pageref*{sec:cobordisms}页}\\
\(d_{\infty}\) &  & 545 & \hyperref[sec:appA-transversality]{第\pageref*{sec:appA-transversality}页}\\
\((D_1f)\) &  & 20 & \hyperref[sec:ch01-exercises]{第\pageref*{sec:ch01-exercises}页}\\
\((D_2f)\) &  & 20 & \hyperref[sec:ch01-exercises]{第\pageref*{sec:ch01-exercises}页}\\
\(E(u)\) & 环路 u 的能量 & 165 & \hyperref[sec:space-of-solutions]{第\pageref*{sec:space-of-solutions}页}\\
\(\widetilde{\exp}\) &  & 481 & \hyperref[sec:operator-frho]{第\pageref*{sec:operator-frho}页}\\
\(\mathscr{F}\) & 弗洛尔映射 & 221 & \hyperref[sec:results]{第\pageref*{sec:results}页}\\
\(\mathscr{F}^{\Gamma}\) &  & 388 & \hyperref[sec:morphism-phi-gamma]{第\pageref*{sec:morphism-phi-gamma}页}\\
\(\mathscr{F}_{\rho}\) &  & 317 & \hyperref[sec:construction-psi-rho]{第\pageref*{sec:construction-psi-rho}页}\\
\(\Phi^{\Gamma}\) &  & 383 & \hyperref[mtfh:chapter11-introduction]{第\pageref*{mtfh:chapter11-introduction}页}\\
\(\mathrm{grad}_x f\) & f 在 x 处的梯度 & 23 & \hyperref[sec:gradients-pseudo-charts]{第\pageref*{sec:gradients-pseudo-charts}页}\\
\(\mathscr{H}\) &  & 272 & \hyperref[sec:fredholm-property]{第\pageref*{sec:fredholm-property}页}\\
\(\mathscr{H}_{\mathrm{reg}}\) &  & 223 & \hyperref[sec:results]{第\pageref*{sec:results}页}\\
\((H^a,J^a)\) &  & 383 & \hyperref[mtfh:chapter11-introduction]{第\pageref*{mtfh:chapter11-introduction}页}\\
\(\mathrm{Ham}(W)\) &  & 517 & \hyperref[sec:ch14-exercises-ch5]{第\pageref*{sec:ch14-exercises-ch5}页}\\
\(\mathrm{Hess}\) & 海塞矩阵 & 138 & \hyperref[sec:hamiltonian-vector-fields]{第\pageref*{sec:hamiltonian-vector-fields}页}\\
\(HM_\star\) & 莫尔斯同调 & 79 & \hyperref[sec:homology]{第\pageref*{sec:homology}页}\\
\(HM_\star(\cdot;\mathbf{Z})\) &  & 79 & \hyperref[sec:homology]{第\pageref*{sec:homology}页}\\
\(HM_\star(\cdot;\mathbf{Z}/2)\) &  & 79 & \hyperref[sec:homology]{第\pageref*{sec:homology}页}\\
\(HM_k\) &  & 6 & \hyperref[mtfh:part1-introduction]{第\pageref*{mtfh:part1-introduction}页}\\
\(H_k(f,X)\) &  & 53 & \hyperref[sec:definition-of-the-complex]{第\pageref*{sec:definition-of-the-complex}页}\\
\(\mathrm{Ind}\) & 指标 & 28 & \hyperref[sec:gradients-pseudo-charts]{第\pageref*{sec:gradients-pseudo-charts}页}\\
\(\mathscr{J}_n\) & 实 2n 维空间上的复结构 & 140 & \hyperref[sec:complex-structures]{第\pageref*{sec:complex-structures}页}\\
\(\mathscr{J}_c(\omega)\) & 由辛形式标定的近复结构 & 142 & \hyperref[sec:complex-structures]{第\pageref*{sec:complex-structures}页}\\
\((\mathscr{J}\times\mathscr{H})_{\mathrm{reg}}\) &  & 223 & \hyperref[sec:results]{第\pageref*{sec:results}页}\\
\(\lambda\) & 刘维尔形式 & 131 & \hyperref[sec:symplectic-manifolds-definition]{第\pageref*{sec:symplectic-manifolds-definition}页}\\
\(\mathcal{L}(\cdot,\cdot),\ \mathscr{L}(\cdot,\cdot)\) & 莫尔斯／弗洛尔部分的记号 & 39，306 & \hyperref[sec:the-smale-condition]{第\pageref*{sec:the-smale-condition}页}\\
\(\mathcal{L}(a,b)\) & 轨道空间 & 39 & \hyperref[sec:the-smale-condition]{第\pageref*{sec:the-smale-condition}页}\\
\(\mathscr{L}(x,y)\) & 轨道空间 & 306 & \hyperref[sec:trajectory-spaces]{第\pageref*{sec:trajectory-spaces}页}\\
\(\overline{\mathcal{L}}(\cdot,\cdot),\ \overline{\mathscr{L}}(\cdot,\cdot)\) & 莫尔斯／弗洛尔部分的记号 & 57，308 & \hyperref[sec:definition-of-the-complex]{第\pageref*{sec:definition-of-the-complex}页}\\
\(\overline{\mathcal{L}}(a,b)\) & 破碎轨迹空间 & 57 & \hyperref[sec:definition-of-the-complex]{第\pageref*{sec:definition-of-the-complex}页}\\
\(\overline{\mathscr{L}}(x,y)\) & 破碎轨迹空间 & 308 & \hyperref[sec:trajectory-spaces]{第\pageref*{sec:trajectory-spaces}页}\\
\(\mathscr{L}(\mathscr{H})\) &  & 275 & \hyperref[sec:fredholm-property]{第\pageref*{sec:fredholm-property}页}\\
\(\mathscr{L}W\) & 环路空间 & 156 & \hyperref[sec:outline-floer-homology]{第\pageref*{sec:outline-floer-homology}页}\\
\(\widetilde{\mathscr{L}W}\) &  & 185 & \hyperref[sec:functions-forms-covers]{第\pageref*{sec:functions-forms-covers}页}\\
\(\Lambda_n\) & 拉格朗日子空间的格拉斯曼流形 & 528 & \hyperref[sec:ch14-exercises-ch7]{第\pageref*{sec:ch14-exercises-ch7}页}\\
\(L_{\rho}\) &  & 318 & \hyperref[sec:construction-psi-rho]{第\pageref*{sec:construction-psi-rho}页}\\
\(\mu\) & 马斯洛夫指标 & 197 & \hyperref[sec:maslov-index-path]{第\pageref*{sec:maslov-index-path}页}\\
\(\mathscr{M}\) &  & 362 & \hyperref[sec:floer-to-morse-results]{第\pageref*{sec:floer-to-morse-results}页}\\
\(\mathcal{M}(\cdot,\cdot),\ \mathscr{M}(\cdot,\cdot)\) & 莫尔斯／弗洛尔部分的记号 & 39，175 & \hyperref[sec:the-smale-condition]{第\pageref*{sec:the-smale-condition}页}\\
\(\mathcal{M}(a,b)\) &  & 39 & \hyperref[sec:the-smale-condition]{第\pageref*{sec:the-smale-condition}页}\\
\(\mathscr{M}(x,y)\) &  & 175 & \hyperref[sec:compactness-proof]{第\pageref*{sec:compactness-proof}页}\\
\(\mathscr{M}(H,J)\) &  & 223 & \hyperref[sec:results]{第\pageref*{sec:results}页}\\
\(\mathscr{M}^{\Gamma}\) &  & 385 & \hyperref[sec:morphism-phi-gamma]{第\pageref*{sec:morphism-phi-gamma}页}\\
\(n_X(a,b)\) &  & 54 & \hyperref[sec:definition-of-the-complex]{第\pageref*{sec:definition-of-the-complex}页}\\
\(N_X(a,b)\) &  & 68 & \hyperref[sec:orientation-complex-over-z]{第\pageref*{sec:orientation-complex-over-z}页}\\
\(\pi_1\) & 基本群 & 90 & \hyperref[sec:homology-connectedness]{第\pageref*{sec:homology-connectedness}页}\\
\(\pi_2\) & 第二同伦群 & 156 & \hyperref[sec:outline-floer-homology]{第\pageref*{sec:outline-floer-homology}页}\\
\(P_V(t)\) & V 的庞加莱多项式 & 86 & \hyperref[sec:euler-characteristic]{第\pageref*{sec:euler-characteristic}页}\\
\(\mathbf{P}^2(\mathbf{R})\) &  & 19 & \hyperref[sec:ch01-exercises]{第\pageref*{sec:ch01-exercises}页}\\
\(\mathbf{P}^n(\mathbf{R})\) &  & 103 & \hyperref[sec:applications]{第\pageref*{sec:applications}页}\\
\(\mathbf{P}^n(\mathbf{C})\) &  & 18 & \hyperref[sec:examples-of-morse-functions]{第\pageref*{sec:examples-of-morse-functions}页}\\
\(\mathscr{P}^{1,p}(x,y)\) &  & 225 & \hyperref[sec:results]{第\pageref*{sec:results}页}\\
\(R(u)\) &  & 247 & \hyperref[sec:transversality]{第\pageref*{sec:transversality}页}\\
\(\mathscr{S}\) &  & 191 & \hyperref[sec:toward-definition-index]{第\pageref*{sec:toward-definition-index}页}\\
\(S^2\) &  & 3 & \hyperref[mtfh:part1-introduction]{第\pageref*{mtfh:part1-introduction}页}\\
\(S^n\) &  & 56 & \hyperref[sec:definition-of-the-complex]{第\pageref*{sec:definition-of-the-complex}页}\\
\(S_k\) &  & 201 & \hyperref[sec:maslov-index-path]{第\pageref*{sec:maslov-index-path}页}\\
\(\mathrm{Sp}(2n)\) & 辛群 & 144 & \hyperref[sec:complex-structures]{第\pageref*{sec:complex-structures}页}\\
\(\mathrm{Sp}(2n)^\star\) &  & 191 & \hyperref[sec:toward-definition-index]{第\pageref*{sec:toward-definition-index}页}\\
\(\mathrm{Sp}(2n)^+\) &  & 194 & \hyperref[sec:toward-definition-index]{第\pageref*{sec:toward-definition-index}页}\\
\(\mathrm{Sp}(2n)^-\) &  & 194 & \hyperref[sec:toward-definition-index]{第\pageref*{sec:toward-definition-index}页}\\
\(T^2\) &  & 17 & \hyperref[sec:examples-of-morse-functions]{第\pageref*{sec:examples-of-morse-functions}页}\\
\(T^\star V\) & V 的余切丛 & 18 & \hyperref[sec:examples-of-morse-functions]{第\pageref*{sec:examples-of-morse-functions}页}\\
\(\pitchfork\) & 横截于 & 38 & \hyperref[sec:the-smale-condition]{第\pageref*{sec:the-smale-condition}页}\\
\(w_{\rho}\) & 预粘合 & 313 & \hyperref[sec:broken-gluing]{第\pageref*{sec:broken-gluing}页}\\
\(W_{\rho}^{\perp}\) &  & 320 & \hyperref[sec:construction-psi-rho]{第\pageref*{sec:construction-psi-rho}页}\\
\(W^{s,p}\) &  & 574 & \hyperref[sec:sobolev-spaces-rn]{第\pageref*{sec:sobolev-spaces-rn}页}\\
\(W^{1,p}\) &  & 574 & \hyperref[sec:sobolev-spaces-rn]{第\pageref*{sec:sobolev-spaces-rn}页}\\
\(W^{1,p}(\mathbf{R}\times S^1;W)\) &  & 223 & \hyperref[sec:results]{第\pageref*{sec:results}页}\\
\(W^{1,p}(x^\star TW)\) &  & 186 & \hyperref[sec:banach-manifold-lw]{第\pageref*{sec:banach-manifold-lw}页}\\
\(\mathcal{W}(\lambda,U^-,U^+)\) &  & 58 & \hyperref[sec:space-of-connections]{第\pageref*{sec:space-of-connections}页}\\
\(W^s(\cdot)\) & 稳定流形 & 27 & \hyperref[sec:gradients-pseudo-charts]{第\pageref*{sec:gradients-pseudo-charts}页}\\
\(W^u(\cdot)\) & 不稳定流形 & 27 & \hyperref[sec:gradients-pseudo-charts]{第\pageref*{sec:gradients-pseudo-charts}页}\\
\(\overline{W^u}(\cdot)\) &  & 112 & \hyperref[sec:cellular-homology]{第\pageref*{sec:cellular-homology}页}\\
\(X_H\) & H 的哈密顿向量场 & 134 & \hyperref[sec:examples-symplectic-manifolds]{第\pageref*{sec:examples-symplectic-manifolds}页}\\
\(X_t\) & 表示随时间变化的哈密顿向量场 & 136 & \hyperref[sec:hamiltonian-vector-fields]{第\pageref*{sec:hamiltonian-vector-fields}页}\\
\(\mathscr{X}_H\) &  & 163 & \hyperref[sec:gradient-floer-equation]{第\pageref*{sec:gradient-floer-equation}页}\\
\(Y\mathbin{\#_{\rho}}Z\) &  & 314 & \hyperref[sec:pregluing]{第\pageref*{sec:pregluing}页}\\
\(Z_i^{\rho}\) &  & 316 & \hyperref[sec:construction-psi-rho]{第\pageref*{sec:construction-psi-rho}页}\\
\(\mathscr{Z}(x,y,J)\) &  & 223 & \hyperref[sec:results]{第\pageref*{sec:results}页}\\
\(\omega\) & 辛形式 & 129 & \hyperref[sec:symplectic-vector-spaces]{第\pageref*{sec:symplectic-vector-spaces}页}\\
\bottomrule
\end{longtable}
}
